\section{Ограничения и бъдещо развитие}
\label{sec:future_work}

Текущата версия на библиотеката вече демонстрира работеща архитектура за \ORMName{} върху релационни и нерелационни backend-и, но научната коректност изисква ясно разграничение между вече постигнатото и следващите стъпки. Библиотеката съдържа реален слой за query IR, capability-базиран конекторен contract, migration прототип, helper-и за нишкова безопасност и набор от unit и integration тестове. Въпреки това част от подсистемите все още трябва да бъдат разширени или валидирани по-задълбочено в production-like условия.

Бъдещото развитие не е произволен списък от идеи. То следва непосредствено от архитектурните решения, анализирани в предходните глави. Ако C++ моделът е първичен носител на логическата схема, следваща естествена стъпка е по-пълна автоматизация на schema-aware tooling. Ако connector слоят е централна точка на extensibility, тогава бъдещата работа трябва да уточни по-фини capability класове. Ако query IR-ът е реално споделен слой между различни backend-и, тогава той трябва да бъде изпитан при по-богати operational сценарии и performance ограничения.

\subsection{Стратегически направления за развитие}

От гледна точка на настоящата дипломна работа бъдещото развитие може да бъде групирано в четири стратегически направления. Първото е \emph{по-голяма operational зрялост} на вече наличните backend-и. Второто е \emph{по-пълна автоматизация на схемата и метаданните}. Третото е \emph{намаляване на runtime overhead-а} чрез по-ниско ниво на изпълнение и по-добър контрол върху wire protocol слоя. Четвъртото е \emph{намаляване на boilerplate-а} чрез по-дълбока интеграция с новите езикови механизми на C++.

\begin{figure}[H]
\centering
\begin{tikzpicture}[node distance=1.35cm and 1.2cm]
    \node[ormaccent] (core) {\textbf{Текуща библиотека}\\typed IR + connector contract + tests};
    \node[ormblock, above right=of core] (maturity) {Operational\\maturity};
    \node[ormblock, below right=of core] (schema) {Schema-aware\\automation};
    \node[ormblock, below left=of core] (perf) {Execution / wire\\optimisation};
    \node[ormblock, above left=of core] (lang) {Language-level\\automation};

    \draw[ormline] (core) -- (maturity);
    \draw[ormline] (core) -- (schema);
    \draw[ormline] (core) -- (perf);
    \draw[ormline] (core) -- (lang);
\end{tikzpicture}
\caption{Четири стратегически направления за бъдещо развитие}
\label{fig:future_work_axes}
\end{figure}

Фигура~\ref{fig:future_work_axes} е полезна, защото показва, че бъдещата работа не е еднопластова. Част от задачите са насочени към експлоатационна зрялост, други --- към по-добра compile-time автоматизация, а трети --- към performance и runtime инженеринг. Това разграничение е важно, защото различните направления изискват различни критерии за успех.

\subsection{Разширяване на live backend покритието}

Макар хранилището вече да съдържа live integration тестове за няколко backend-а, зрелостта на отделните реализации не е еднаква. SQLite остава най-естественият референтен релационен път, защото съчетава най-плътно свързани query rendering, bind логика и result hydration. При PostgreSQL, MySQL, MongoDB, Redis, Neo4j и Cassandra следващата стъпка е по-широко покриване на operational сценарии: по-богати типове, диагностични пътища при грешка, lifecycle политики за връзки и оценка на поведението при по-сложни заявки.

Това направление е важно не само като инженерно подобрение, а и като научна проверка на границите на общия query IR. Само при по-широко емпирично покритие може да се оцени доколко един и същ логически слой остава адекватен за различни типове системи. Например join-базирана логика, която е естествена за PostgreSQL, има съвсем различен operational смисъл в Neo4j или Cassandra. Следователно live покритието не е просто тестова дисциплина, а средство за изследване на архитектурната приложимост.

Допълнително, live backend покритието трябва да обхване и негативните пътища: грешки при connection establishment, частично невалидни параметри, driver-specific timeout сценарии и поведение при прекъсване на връзката. Тези въпроси почти винаги са подценени в ранните ORM библиотеки, но в multi-backend контекст те са решаващи за реалната употреба.

\subsection{Пълна schema-aware интеграция и runtime introspection}

Прототипният migration слой е достатъчен, за да покаже как C++ entity моделът може да участва в описване на schema evolution. Следващото естествено направление е live introspection на реалната схема, извличане на metadata от драйверите и безопасно изпълнение на DDL операции. Това е особено важно за релационните backend-и, но не е без значение и за документни или wide-column системи, където понятието за schema може да бъде по-гъвкаво или частично.

По-нататъшната работа в тази посока трябва да изясни и границата между compile-time гаранции и runtime факти за вече разположена система. Компилаторът може да валидира query структурата, но не и реалното състояние на production schema без допълнителна runtime информация. Следователно бъдещият migration слой трябва да комбинира два източника на истина: C++ entity описанието и извлечената metadata картина на реалната база.

\begin{figure}[H]
\centering
\begin{tikzpicture}[node distance=1.3cm and 1.1cm]
    \node[ormaccent] (model) {C++ model\\entity metadata};
    \node[ormblock, right=of model] (live) {Live backend\\metadata};
    \node[ormblock, below=of $(model)!0.5!(live)$] (compare) {Schema diff +\\compatibility analysis};
    \node[ormblock, below left=of compare] (ddl) {DDL generation};
    \node[ormblock, below right=of compare] (report) {Safety report\\and dry-run};
    \node[ormblock, below=of $(ddl)!0.5!(report)$] (apply) {Controlled migration\\execution};

    \draw[ormline] (model) -- (compare);
    \draw[ormline] (live) -- (compare);
    \draw[ormline] (compare) -- (ddl);
    \draw[ormline] (compare) -- (report);
    \draw[ormline] (ddl) -- (apply);
    \draw[ormline] (report) -- (apply);
\end{tikzpicture}
\caption{Целева посока за schema-aware tooling отвъд текущия migration прототип}
\label{fig:future_schema_tooling}
\end{figure}

Фигура~\ref{fig:future_schema_tooling} показва, че бъдещият migration pipeline трябва да бъде двупосочен: едновременно да чете от compile-time описанието и от живата система. Това е естествена следваща стъпка, ако библиотеката трябва да бъде използвана не само експериментално, но и като practical development tool.

\subsection{Нишкова безопасност и жизнен цикъл на връзките}

Съществуващите helper-и за connection pool, thread-local достъп и transaction guard показват, че библиотеката вече съдържа архитектурна основа за multi-threaded употреба. Бъдещото развитие трябва да отговори на по-конкретни въпроси: как се дефинира \code{supports\_concurrent\_execute} за различните драйвери; кога една връзка е безопасна за споделяне; какви са правилата за реконекция, затваряне и отложено освобождаване на ресурси.

Тук има и важен изследователски аспект: различните backend-и имат различна concurrency семантика, а следователно compile-time capability обозначаването трябва да бъде достатъчно прецизно, за да не създава подвеждащи обещания. Напълно възможно е бъдещи версии на библиотеката да разделят общото \code{supports\_concurrent\_execute} на по-фини capability класове като безопасно паралелно четене, паралелно изпълнение с отделни connection handle-и или строга serialization политика.

Практически това означава, че thread-safety helper слоят не трябва да остане изолиран помощен модул. Напротив, той трябва да се превърне в архитектурно продължение на connector contract-а и да участва в пълната оценка на зрелостта на даден backend.

\subsection{Ниско ниво на изпълнение и wire protocol оптимизации}

Експерименталният слой в \path{WireProtocol/wire_protocol.hpp} показва възможни бъдещи посоки като \code{io\_uring} / IOCP awaitable-и, \code{zero\_copy\_result}, \code{batch\_insert} и \code{make\_constexpr\_sql}. Тези елементи все още не са production-ready комуникационен стек, но очертават възможност библиотеката да се разшири отвъд класическия synchronous driver модел.

Особено перспективни са нулево-копийната обработка на резултати и compile-time генерирането на SQL за конектори, които го позволяват. Те биха приближили библиотеката до още по-нисък runtime overhead, без да се жертва типизираната архитектура. Това е важно, защото една честа критика към богато типизираните abstraction слоеве е, че скриват прекомерен runtime разход. Бъдещата работа трябва да покаже, че compile-time изразителността и ниският overhead не са задължително противоположни цели.

Освен това performance темата не трябва да се разглежда само през latency или throughput. Тя включва и memory layout на резултатите, броя временни allocation-и, повторната употреба на prepared query артефакти и качеството на driver-level batching стратегиите.

\subsection{Интеграция с C++26 reflection и допълнителна автоматизация}

Поддръжката на C++26 reflection е вече концептуално подготвена чрез detection логика и helper-и, а \path{tests/unit/test_cpp26_reflection.cpp} показва наличието на early verification path. Публичният API обаче все още използва експлицитни string literal имена за \code{property<T, Name>}. Бъдещото развитие трябва да превърне reflection path-а в стабилен публичен механизъм, който позволява автоматично извличане на имената на полетата, без това да нарушава обратната съвместимост.

Подобно развитие би намалило boilerplate-а и би направило още по-видима идеята, че C++ моделът е първичният източник на логическата схема. От изследователска гледна точка то би било интересно и защото би доближило библиотеката до model-first стил, при който голяма част от metadata информацията се извлича автоматично, а не се описва повторно.

Тук обаче има и важна граница. Reflection не трябва да разрушава яснотата на compile-time contract-а. Ако автоматизацията намали прозрачността на типа или затрудни съобщенията за грешка, тя може да отслаби една от основните ценности на библиотеката. Следователно тази посока трябва да се развива внимателно.

\subsection{Разширяване на connector contract-а и capability таксономията}

С добавянето на нови backend-и ще стане необходимо по-прецизно разграничаване между различни класове възможности. Възможно е например бъдещи версии на библиотеката да разграничават отделно graph traversal capability, partial document embedding support, schema introspection support, streaming result support или explicit batching support. Тази посока следва естествено от текущата trait-базирана архитектура.

По-нататъшната работа трябва да пази едно важно правило: новите capability-та не бива да бъдат декоративни. Те трябва да носят реално compile-time значение и да ограничават недопустимите операции. Ако capability tag няма архитектурно последствие, той не добавя стойност. Точно обратното --- прави contract-а по-шумен и по-малко точен.

\begin{table}[H]
\centering
\small
\caption{Приоритетни направления за бъдещо развитие}
\label{tab:future_work_priorities}
\begin{tabularx}{\textwidth}{L{3.0cm}L{3.2cm}L{3.1cm}Y}
\toprule
\textbf{Направление} & \textbf{Засегнати подсистеми} & \textbf{Очакван ефект} & \textbf{Основно предизвикателство} \\
\midrule
Live backend maturity & connector-и, integration tests, error paths & по-висока practical reliability & различна operational семантика между backend-ите \\
Schema-aware automation & \code{migrate<DB>}, DDL hooks, metadata introspection & по-близка до production ORM интеграция & съчетаване на compile-time и runtime истина \\
Execution optimisation & wire protocol layer, prepared queries, batching & по-нисък runtime overhead & запазване на typed abstraction-а \\
Reflection-driven modelling & entity layer, metadata extraction, API ergonomics & по-малко boilerplate и по-автоматичен model-first подход & баланс между автоматизация и прозрачност \\
Capability refinement & \code{connector\_trait<DB>} и static checks & по-точен contract за нови backend-и & избягване на декоративни capability tag-ове \\
\bottomrule
\end{tabularx}
\end{table}

\subsection{Необходимост от по-строга емпирична оценка}

Една бъдеща версия на работата следва да включва и по-систематичен benchmarking слой. Той трябва да сравнява не само отделни backend-и, но и различни execution режими вътре в библиотеката --- директно изпълнение, \code{prepare()} + reuse, batch insert сценарии и различни форми на result materialization. Такъв benchmarking би бил ценен както инженерно, така и научно, защото би показал реалната цена на различните архитектурни избори.

Наред с това, полезно би било да се добави и формализирана оценка на error diagnostics. Една от силните страни на compile-time ORM дизайна е именно ранното локализиране на грешки. Следователно качеството на съобщенията за грешка и яснотата на compile-time failure пътищата също са валиден обект на изследване.

\subsection{Обобщение}

Бъдещото развитие на библиотеката не започва от празно място. Напротив, то стъпва върху вече работеща архитектура, която е достатъчно обща, за да поеме нови backend-и, и достатъчно строга, за да наложи формални ограничения. Именно това прави разработката подходяща както за продължаване като инженерна система, така и за по-нататъшни академични изследвания в областта на compile-time моделирането на достъп до данни.

Най-важното заключение от настоящата глава е, че бъдещата работа не трябва да се разбира като поправка на несъстоятелен прототип. Тя е естествено разгръщане на вече защитима архитектурна основа. Това е съществено различие: библиотеката е достигнала ниво, на което разширението е въпрос на дисциплинирано развитие, а не на концептуално преосмисляне.
